\chapter{\texttt{pipeline.earth} --- Earth surface to detector}
\label{ch:pipeline-earth}

\texttt{pipeline.earth} propagates an already-known incident state --
flavour-basis coherent amplitudes, or incoherent mass-basis weights such as
the ones \cref{ch:pipeline-solar} produces -- through the Earth to a
detector, over either an explicit nadir-angle grid or the configured
day/night exposure window (\cref{sec:earth-exposure}). It is the Earth leg
reused, unchanged, by \texttt{pipeline.solar\_earth}
(\cref{ch:pipeline-solar-earth}); used on its own it also covers any
externally-produced incident flux (reactor, accelerator, atmospheric
production already reduced to a single flavour, \ldots) that only needs
Earth regeneration, not a solar production stage.

\section{\texttt{pipeline.earth} --- pure Earth workflow}

\modulehead{pipeline/earth.py}{\pyfunc{propagate\_earth\_to\_detector} and
\pyfunc{propagate\_earth\_to\_detector\_exposure}: Earth-crossing
propagation over an explicit nadir-angle grid, or averaged over the
configured exposure window.}

\begin{implbox}
\begin{itemize}
  \item \pyclass{EarthDetectorResult} (frozen dataclass): \code{incident\_state},
    \code{incident\_basis} (\code{"flavour"} or \code{"mass"}), \code{E\_MeV},
    the built \code{profile} (\pyclass{medium.earth.profile.EarthProfile}),
    and \emph{either} the explicit-grid pair (\code{eta}, \code{exposure},
    \code{probabilities\_eta}) \emph{or} the exposure-averaged
    \code{probabilities\_exposure} -- exactly one of the two pairs is
    populated, matching which of the two functions below built the result.
  \item \pyfunc{propagate\_earth\_to\_detector(incident\_state, *, E\_MeV,
    config, incident\_basis="flavour", earth\_profile=None, eta=None)}:
    resolves the nadir-angle grid via
    \pyfunc{medium.earth.exposure\_table.prepare\_nadir\_exposure} (an
    explicit \code{eta} if given, otherwise the one implied by
    \code{config.exposure}), then calls
    \pyfunc{earth\_probability\_state} (\cref{ch:earth}) on the Cartesian
    product of the energy and eta grids. \code{config.earth.chunk\_eta}
    splits the eta axis into sub-batches evaluated one at a time (a memory
    control for large energy--angle grids), concatenated back together
    before being returned.
  \item \pyfunc{propagate\_earth\_to\_detector\_exposure(incident\_state, *,
    E\_MeV, config, incident\_basis="flavour", earth\_profile=None)}: calls
    \pyfunc{earth\_probability\_exposure} (\cref{sec:earth-exposure})
    directly, i.e.\ propagates and averages over the nadir-angle exposure
    weight $W(\eta)$ in one step rather than returning a
    per-angle grid.
\end{itemize}
Both functions accept \code{incident\_basis="mass"} for an incoherent mass
mixture (e.g.\ \cref{ch:pipeline-solar}'s \code{mass\_weights}) and
\code{"flavour"} (the default) for coherent flavour amplitudes, dispatching
to \pyfunc{earth\_probability\_state}/\pyfunc{earth\_probability\_exposure}'s
own \code{massbasis} flag accordingly.
\end{implbox}

\begin{examplebox}[title={A pure $\nu_\mu$ beam crossing the Earth on a nadir-angle grid}]
\begin{lstlisting}[style=tpython]
import torch

from tpeanuts.config.propagation import PropagationConfig
from tpeanuts.core.common.neutrino import flavour_state
from tpeanuts.pipeline.earth import propagate_earth_to_detector
from tpeanuts.util.context import RuntimeContext

context = RuntimeContext.resolve("cpu", torch.float64)
oscillation = PropagationConfig.oscillation_parameters_from_preset(
    "_SM_NUFIT52_NO", context=context,
)
config = PropagationConfig(runtime=context, oscillation=oscillation)

nu_mu = flavour_state("mu", device=context.device, dtype=context.dtype)

result = propagate_earth_to_detector(
    nu_mu,
    E_MeV=[500.0, 2000.0],
    config=config,
    eta=[0.2, 0.8, 1.4],   # nadir angle in radians, 0 = through the core
)

print(result.probabilities_eta.shape)   # torch.Size([2, 3, 3]): (E, eta, flavour)
\end{lstlisting}
\end{examplebox}

For a detector-exposure average instead of a fixed grid, call
\pyfunc{propagate\_earth\_to\_detector\_exposure} with the same
\code{nu\_mu}/\code{E\_MeV}/\code{config} -- no \code{eta} argument, since
the angular average is taken internally over
\code{config.exposure} (\cref{sec:earth-exposure}):
\begin{examplebox}[title={The same beam, averaged over the configured detector exposure}]
\begin{lstlisting}[style=tpython]
from tpeanuts.pipeline.earth import propagate_earth_to_detector_exposure

exposure_result = propagate_earth_to_detector_exposure(
    nu_mu,
    E_MeV=[500.0, 2000.0],
    config=config,
)
print(exposure_result.probabilities_exposure.shape)   # torch.Size([2, 3]): (E, flavour)
\end{lstlisting}
\end{examplebox}
